\SourcePage{529}{532}
\section{Matrix elements of tensors}

In \S~29 we obtained formulas giving the dependence of the matrix elements of a vector physical quantity on the value of the angular-momentum projection. In fact, those formulas are a special case of general formulas that solve the same problem for an irreducible tensor (see p.~268) of any rank\footnote{The treatment of the questions considered in \S~107--109, and most of the results presented there, are due to Racah (1942/1943).}.

The set of $2k+1$ components of an irreducible tensor of rank $k$ ($k$ is an integer) is equivalent, in its transformation properties, to the set of $2k+1$ spherical functions $Y_{kq}$ ($q=-k,\ldots,k$) (see the note on p.~268). This means that suitable linear combinations of the tensor components can be formed so as to give a set of quantities transforming under rotations in the same way as the functions $Y_{kq}$. We shall call such a set of quantities, denoted here by $f_{kq}$, a \emph{spherical tensor} of rank $k$.

Thus a vector corresponds to $k=1$, and the quantities $f_{1q}$ are related to the vector components by
\begin{equation}
 f_{10}=ia_z,\qquad f_{1,\pm1}=\mp\frac{i}{\sqrt2}(a_x\pm ia_y)
 \tag{107.1}
\end{equation}
(cf. (57.7)). The analogous formulas for a tensor of second rank

\SourcePage{530}{533}
are
\begin{equation}
 \begin{gathered}
 f_{20}=-\sqrt{3/2}\,a_{zz},\qquad f_{2,\pm1}=\pm(a_{xz}\pm ia_{yz}),\\
 f_{2,\pm2}=-\frac12(a_{xx}-a_{yy}\pm2ia_{xy}),
 \end{gathered}
 \tag{107.2}
\end{equation}
where $a_{xx}+a_{yy}+a_{zz}=0$\footnote{It is understood that the quantities $f_{kq}$ are complex only because of the change to spherical coordinates; that is, the original Cartesian tensor components are real.}.

Tensor products of two (or more) spherical tensors $f_{k_1q_1}$, $g_{k_2q_2}$ are formed according to the rules for addition of angular momenta, with $k_1$, $k_2$ formally playing the role of the ``angular momenta'' associated with these tensors. Thus, from two spherical tensors of ranks $k_1$ and $k_2$ one can form spherical tensors of ranks $K=k_1+k_2,\ldots,|k_1-k_2|$ by
\begin{equation}
 \begin{aligned}
 (f_{k_1}g_{k_2})_{KQ}&=\sum_{q_1q_2}(q_1q_2|KQ)f_{k_1q_1}g_{k_2q_2}=\\
 &=(-1)^{k_1-k_2+Q}\sqrt{2K+1}\sum_{q_1q_2}
 \begin{pmatrix}k_1&k_2&K\\q_1&q_2&-Q\end{pmatrix}f_{k_1q_1}g_{k_2q_2}
 \end{aligned}
 \tag{107.3}
\end{equation}
(cf. (106.9)). The scalar product of two spherical tensors of the same rank $k$ is, however, conventionally defined by
\begin{equation}
 (f_kg_k)_{00}=\sum_q(-1)^{k-q}f_{kq}g_{k,-q},
 \tag{107.4}
\end{equation}
which differs by a factor $\sqrt{2k+1}$ from the definition supplied by (107.3) with $K=Q=0$ (cf. (106.2))\footnote{If $\mathbf A$ and $\mathbf B$ are two vectors corresponding, by (107.1), to the spherical tensors $f_{1q}$ and $g_{1q}$, then
\[
 (f_1g_1)_{00}=\mathbf{AB}.
\]}. This definition may also be written
\[
 (f_kg_k)_{00}=\sum_q f_{kq}g_{kq}^{*},
\]
on observing that complex conjugation of a spherical tensor is performed according to
\[
 f_{kq}^{*}=(-1)^{k-q}f_{k,-q}
\]
(cf. (28.9))\footnote{We repeat here the observation made above in connection with (106.8): with this rule, complex conjugation of the tensors of ranks $k_1$ and $k_2$ on the right-hand side of (107.3) produces the corresponding conjugation of the tensor of rank $K$ on its left-hand side.}.

\SourcePage{531}{534}
Representing physical quantities as spherical tensors is especially convenient when their matrix elements are calculated, since the results of the theory of angular-momentum addition can then be used.

By the definition of matrix elements,
\begin{equation}
 \widehat f_{kq}\psi_{njm}=\sum_{n'j'm'}\langle n'j'm'|f_{kq}|njm\rangle\psi_{n'j'm'},
 \tag{107.5}
\end{equation}
where $\psi_{njm}$ are the wave functions of stationary states of the system, specified by the magnitude $j$ of its angular momentum, its projection $m$, and the set $n$ of the remaining quantum numbers. In their transformation properties, the functions on the right- and left-hand sides of (107.5) correspond respectively to the functions on the right- and left-hand sides of (106.11). The selection rules follow at once.

The matrix elements of the components $f_{kq}$ of an irreducible tensor of rank $k$ can be nonzero only for transitions $jm\to j'm'$ satisfying the ``angular-momentum addition rule'' $\mathbf j'=\mathbf j+\mathbf k$; the numbers $j'$, $j$, $k$ must satisfy the ``triangle rule'' (that is, they can be the lengths of the sides of a closed triangle), and $m'=m+q$. In particular, diagonal matrix elements can be nonzero only if $2j\geqslant k$.

The same analogy of transformation properties further shows that the coefficients in the sum (107.5) must be proportional to those in (106.11) (the \emph{Wigner--Eckart theorem}). This determines their dependence on $m$ and $m'$. Accordingly, we write the matrix elements as
\begin{equation}
 \langle n'j'm'|f_{kq}|njm\rangle=i^k(-1)^{j_{\max}-m'}
 \begin{pmatrix}j'&k&j\\-m'&q&m\end{pmatrix}
 \langle n'j'\|f_k\|nj\rangle,
 \tag{107.6}
\end{equation}
where $j_{\max}$ is the larger of $j$ and $j'$, while $\langle n'j'\|f_k\|nj\rangle$ are quantities independent of $m$, $m'$, and $q'$; they are called \emph{reduced matrix elements}. This formula answers the question of determining the dependence of matrix elements on angular-momentum projections. That dependence is entirely governed by symmetry with respect to the rotation group, whereas dependence on the other quantum numbers is determined by the physical nature of the quantities $f_{kq}$ themselves\footnote{Among the consequences of these results are the selection rules for vector matrix elements stated in \S~29 and the corresponding formulas (29.7)--(29.9).}.

\SourcePage{532}{535}
The operators $\widehat f_{kq}$ are related by
\begin{equation}
 \widehat f_{kq}^{+}=(-1)^{k-q}\widehat f_{k,-q}.
 \tag{107.7}
\end{equation}
Their matrix elements therefore satisfy
\begin{equation}
 \langle n'j'm'|f_{kq}|njm\rangle^{*}=(-1)^{k-q}\langle njm|f_{k,-q}|n'j'm'\rangle.
 \tag{107.8}
\end{equation}
Substitution of (107.6), together with the properties (106.5), (106.6) of the $3j$ symbols, gives the following ``Hermiticity'' relation for the reduced matrix elements\footnote{The phase factor in the definition (107.6) was chosen precisely so that this equality holds.}:
\begin{equation}
 \langle n'j'\|f_k\|nj\rangle=\langle nj\|f_k\|n'j'\rangle^{*}.
 \tag{107.9}
\end{equation}

The matrix elements of the scalar (107.4) are diagonal in $j$ and $m$. The rule for matrix multiplication gives
\[
 \begin{aligned}
 \langle n'jm|(f_kg_k)_{00}|njm\rangle
 ={}&\sum_q(-1)^{k-q}\sum_{n''j''m''}
 \langle n'jm|f_{kq}|n''j''m''\rangle\\
 &\times\langle n''j''m''|g_{k,-q}|njm\rangle.
 \end{aligned}
\]
Substituting (107.6) here and summing over $q$ and $m''$ by means of the orthogonality relation (106.12) for the $3j$ symbols, we obtain
\begin{equation}
 \langle n'jm|(f_kg_k)_{00}|njm\rangle=
 \frac1{2j+1}\sum_{n''j''}\langle n'j\|f_k\|n''j''\rangle
 \langle n''j''\|g_k\|nj\rangle.
 \tag{107.10}
\end{equation}

In the same way one readily obtains the following formulas for sums of squared matrix elements:
\begin{equation}
 \sum_{qm'}|\langle n'j'm'|f_{kq}|njm\rangle|^2=
 \frac1{2j+1}|\langle n'j'\|f_k\|nj\rangle|^2,
 \tag{107.11}
\end{equation}
\begin{equation}
 \sum_{mm'}|\langle n'j'm'|f_{kq}|njm\rangle|^2=
 \frac1{2k+1}|\langle n'j'\|f_k\|nj\rangle|^2.
 \tag{107.12}
\end{equation}
In the first, the sum is over $q$ and $m'$ at fixed $m$; in the second, it is over $m$ and $m'$ at fixed $q$ (with $m'=m+q$ throughout).

For reference, consider the case in which the quantities $f_{kq}$ are the spherical functions $Y_{kq}$ themselves, and give their matrix elements for transitions between states

\SourcePage{533}{536}
of one particle with integral orbital angular momenta $l_1$ and $l_2$, that is, the integrals
\begin{equation}
 \langle l_1m_1|Y_{lm}|l_2m_2\rangle=\int Y_{l_1m_1}^{*}Y_{lm}Y_{l_2m_2}\,do.
 \tag{107.13}
\end{equation}
Besides the selection rule arising from angular-momentum addition ($\mathbf l+\mathbf l_2=\mathbf l_1$), these matrix elements obey the additional rule that $l+l_1+l_2$ must be even. This is connected with parity conservation: the product $(-1)^{l_1+l_2}$ of the parities of the two states must equal the parity $(-1)^l$ of the physical quantity under consideration (see \S~30).

The matrix elements (107.13) are a special case of a more general integral that will be evaluated in \S~110 (see the note on p.~547). They are
\begin{equation}
 \begin{aligned}
 \langle l_1m_1|Y_{lm}|l_2m_2\rangle={}&(-1)^{m_1}i^{-l_1+l_2+l}
 \begin{pmatrix}l_1&l&l_2\\-m_1&m&m_2\end{pmatrix}
 \begin{pmatrix}l_1&l&l_2\\0&0&0\end{pmatrix}\\
 &\times\left[\frac{(2l+1)(2l_1+1)(2l_2+1)}{4\pi}\right]^{1/2}.
 \end{aligned}
 \tag{107.14}
\end{equation}
In particular, for $m_1=m_2=m=0$ we find the following value of the integral of the product of three Legendre polynomials:
\begin{equation}
 \int_{-1}^{1}P_l(\mu)P_{l_1}(\mu)P_{l_2}(\mu)\,d\mu=
 2\begin{pmatrix}l_1&l&l_2\\0&0&0\end{pmatrix}^{2}.
 \tag{107.15}
\end{equation}
